




% The GNN encoder $f_i:
% % \mathbb{R}^{n\timesn} \times \mathbb{R}^{n\times|d|} \rightarrow \mathbb{R}^{n\times|q|}$ 
% \mathbb{R}^{n\times n} \times \mathbb{R}^{n \times q} \rightarrow \mathbb{R}^{n\times d}$ 
% extracts hidden node features $\mathbf{H}_i\in \mathbb{R}^{n\times d}$ from the $i$-th augmented graph $(\tilde{\mathbf{A}}_i, \tilde{\mathbf{X}}_i)$. 
% %
% Usually, multiple encoders are applied to obtain the hidden node features $\mathbf{H}_i$ of different views as:
% \begin{equation}
%     \mathbf{H}_1 = f_1 (\tilde{\mathbf{A}}_1, \tilde{\mathbf{X}}_1), \cdots, \mathbf{H}_k = f_k (\tilde{\mathbf{A}}_k, \tilde{\mathbf{X}}_k).
% \end{equation}
% The GNN encoders are implemented as a 2-layer Graph Convolution Network (GCN):
% \vspace{-0.3cm}
%     % \begin{equation}
%     % \begin{aligned}
%     % \mathrm{GCN}_{l} (\mathbf{X}, \mathbf{A}) &=\sigma\left (\hat{\mathbf{D}}^{-\frac{1}{2}} \hat{\mathbf{A}} \hat{\mathbf{D}}^{-\frac{1}{2}} \mathbf{X} \mathbf{W}\right), \\
%     % f (\mathbf{X}, \mathbf{A}) &=\mathrm{GCN}_{2}\left (\mathrm{GCN}_{1} (\mathbf{X}, \mathbf{A}), \mathbf{A}\right),
%     % \end{aligned}
%     % \end{equation}
%     \begin{equation}
%     \begin{aligned}
%     \mathrm{GCN}_{l} (\mathbf{X}, \mathbf{A}) &=\sigma\left (\hat{\mathbf{D}}^{-\frac{1}{2}} \hat{\mathbf{A}} \hat{\mathbf{D}}^{-\frac{1}{2}} \mathbf{X} \mathbf{W}_l\right), \\
%     f (\mathbf{X}, \mathbf{A}) &=\mathrm{GCN}_{2}\left (\mathrm{GCN}_{1} (\mathbf{X}, \mathbf{A}), \mathbf{A}\right),
%     \end{aligned}
%     \end{equation}
% where $\hat{\mathbf{A}}=\mathbf{A}+\mathbf{I}$ is the adjacency matrix with self-loops, $\mathbf{D}$ is the degree matrix, $\sigma (\cdot)$ is an activation fun. \eg, $\operatorname{ReLU} (\cdot)=\max (0, \cdot)$, and $\mathbf{W}_l$ is a trainable weight matrix for the $l$-th layer of GCNs.
% % \item 

% \vspace{0.1cm}

% \noindent\textbf{Projection Head.} The projection head, $\theta(\cdot)$ is a small network that maps representations to the space where contrastive loss is applied. It is implemented as a multi-layer perceptron (MLP) with one hidden layer to obtain $\mathbf{Z}_{i}=g\left (\mathbf{H}_{i}\right)=\mathbf{W}^{ (2)} \sigma\left (\mathbf{W}^{ (1)} \mathbf{H}_{i}\right)$, where $\sigma$ is the ReLU non-linearity. 
% % \item 

% \vspace{0.1cm}
% \noindent\textbf{Contrastive Loss}. Given two feature matrices $\mathbf{U} \in \mathbb{R}^{n \times d}$ and $\mathbf{V} \in \mathbb{R}^{n\times d}$, where $\mathbf{U}=\mathbf{Z}_1$ and $\mathbf{V} = \mathbf{Z}_2$ are node features obtained from two different views. For any node $i$, its embedding generated in one view, $\mathbf{U}_{i:}$, is treated as the anchor, its embedding generated in the other view, $\mathbf{V}_{i:}$, forms the positive sample. Remaining embeddings of nodes $\mathbf{U}_{j:}$ and $\mathbf{V}_{j:}$ such that $i\neq j$ given the two views are naturally regarded as negative samples. The contrastive loss function $\mathcal{L}$ for all positive pairs is defined as:
% % \begin{equation}
% % \label{eq:mvcontrastiveloss}
% % \begin{aligned}
% % &\mathcal{L} = \frac{1}{n}\sum_{i=1}^{n} \left[{\theta (\mathbf{U}_{i:}, \mathbf{V}_{i:})}/\tau -\log \left(
% % \sum_{j=1}^{n} e^{\theta\left (\mathbf{U}_{i:}, \mathbf{V}_{j:}\right) / \tau}+\sum_{j=1}^{n} e^{\theta\left (\mathbf{U}_{i:}, \mathbf{U}_{j:}\right) / \tau}\right) \right].
% % \end{aligned}
% % \end{equation}

% Note that computing Eq.~(\ref{eq:mvcontrastiveloss}) is both memory costly and time consuming as it requires the computation of three large similarity matrices, $\mathbf{UV}^\top$, $\mathbf{UU}^\top$, $\mathbf{VU}^\top \in \mathbb{R}^{n\times n}$ for two views. Therefore, the memory consumption and runtime  scale %linearly 
% \wrt the product of the number of views and the number of nodes, making such a multi-view setting challenging to run given large-scale graphs.